Pour chacune des affirmations suivantes, indiquer si elle est vraie ou fausse.

Chaque réponse doit être justifiée. Une réponse non justifiée ne rapporte aucun point.

Les affirmations sont indépendantes.

\medskip

\begin{enumerate}
	\item On considère l'équation différentielle \quad $(E)\quad y' - 2y = (x- 3)\e^x$,
	\quad
	où $y$ est une fonction de la variable $x$, définie et dérivable sur $\R$.

	\textbf{Affirmation 1} :

	La fonction $f$ définie sur $\R$ par : $f(x) = (x - 2)\e^x$ est solution de l'équation différentielle (E).

	\item Soit $f$ la fonction définie sur $\R$, par \quad $f(x) = (2x + 1) \cos (x)$.

	On intègre par parties $\displaystyle\int_0^{\pi} f(x)\:\text{d}x$.

	\textbf{Affirmation 2} :
	\[\displaystyle\int_0^{\pi} f(x)\:\text{d}x =  - 4\pi\]

	\item On considère l'inéquation suivante sur $]7 ; +\infty[$ :
	\[\ln \left(2x^2 - 6x - 56\right) \geqslant \ln(4x - 8)\]

	\textbf{Affirmation 3} :

	L'ensemble des solutions de cette inéquation est l'intervalle $[8 ; +\infty[$.

	\item On considère la suite $(u_n)$ définie par $u_0 = 1$ et, pour tout entier naturel $n$,
	\[u_{n+1} = \dfrac{ 3u_n}{5u_n + 3}\]

	On admet que $u_n \ne 0$ pour tout entier naturel $n$.

	On considère la suite $(v_n)$ définie pour tout entier naturel $n$ par \quad $v_n = \dfrac{2}{u_n}$.

	\textbf{Affirmation 4} :

	La suite $(v_n)$ est une suite arithmétique.
	\item On lance \np{3600} fois un dé équilibré à six faces numérotées de 1 à 6.

	\textbf{Affirmation 5} :

	La probabilité que le nombre d'apparitions du 6 soit strictement compris entre $450$ et $750$ est minorée par $0,97$.
\end{enumerate}

\needspace{20\baselineskip}
